\chapter{Prompt Engineering Artifacts}
\label{app:prompts}

The platform delegates four tasks to Large Language Models: parsing a CV into
structured data, extracting hiring requirements from a job description, and
conducting the two modes of the AI Interviewer (technical-skill and language).
Each task is driven by a hand-written, schema-constrained system prompt. This
appendix reproduces the core of those prompts
(\secref{app:prompts-application}) and, separately, documents how AI coding
assistants were used as a bounded tool during development
(\secref{app:prompts-development}).

Two conventions apply throughout. First, the prompts are reproduced from the
source files (\texttt{src/server/infrastructure/ai/} for the Next.js services
and \texttt{ai\_interviewer\_backend/} for the FastAPI sidecar); long enumerated
value lists and repeated few-shot examples are elided with \texttt{[\ldots]}
purely to keep each listing to one page. Second, every prompt ends by
constraining the model to emit JSON only, which the calling code then validates
with Zod (Next.js) or Pydantic (FastAPI) before it is trusted.

%----------------------------------------------------------------------------------------

\section{Application Prompts}
\label{app:prompts-application}

\subsection{CV Parsing}
\label{app:prompts-cv}

The CV parser converts the raw text of an uploaded CV into the candidate
schema. It is called once per CV with \texttt{temperature = 0.1} and
\texttt{response\_format = json\_object}. The prompt fixes the output shape,
pins the canonical field-of-work vocabulary, and asks the model to self-report
a per-field confidence so that low-confidence extractions can be flagged for
review rather than trusted silently.

\begin{lstlisting}[caption={CV parsing system prompt (abridged).},label={lst:prompt-cv}]
You are a CV/resume parser. Extract structured information from the
following CV text.

Return a JSON object with exactly this structure:
{
  "firstName": "string",
  "lastName": "string",
  "email": "string or null",
  "location": "city, country or null",
  "experiences": [
    { "jobTitle": "string", "company": "string or null",
      "startDate": "YYYY-MM", "endDate": "YYYY-MM or null",
      "isCurrent": false, "fieldsOfWork": ["Retail", "Sales"] }
  ],
  "education": [
    { "institution": "string or null", "degree": "string or null",
      "level": "HIGH_SCHOOL | BACHELOR | MASTER | PHD | VOCATIONAL | OTHER" }
  ],
  "languages": [ { "language": "name", "level": "A1..C2 or null" } ],
  "skills":    [ { "name": "skill", "category": "Technical | Soft Skill | Tool | Language | Other" } ],
  "estimatedTotalYears": 5.5,
  "parsingConfidence": {
    "overall": 0.85, "name": 0.99, "location": 0.70,
    "flags": ["location_uncertain", "missing_language_levels", "date_gaps", "no_email"]
  }
}

Official fields of work (use ONLY these exact strings; do NOT invent others):
  - Retail   - Sales   - Data   - Design   - Digital   - Finance
  - Technology   - Supply Chain & Sourcing   - People & Culture   [...]

Rules:
- firstName and lastName are REQUIRED -- never return null; infer from the
  email or headers if needed, and only use "Unknown" as a last resort.
- For all other fields, if not found, use null.
- Map languages to CEFR (native/fluent = C2, advanced = C1, intermediate = B1-B2,
  basic = A1-A2).
- Include not only formal degrees but also certifications, courses and bootcamps,
  classified as VOCATIONAL or OTHER.
- For each experience pick 1-3 fieldsOfWork from the list above; return [] rather
  than invent a value.
- Use the canonical name of each skill ("JavaScript" not "JS scripting") and split
  compound entries ("HTML, CSS & JavaScript" -> three skills).
- For parsingConfidence, honestly assess certainty (0.0-1.0) and add a flag for
  every uncertainty.
- Return ONLY valid JSON, no markdown or extra text.
\end{lstlisting}

\subsection{Job-Description Requirement Extraction}
\label{app:prompts-jd}

The requirements extractor turns free-text job descriptions --- which arrive in
several languages and rarely follow a fixed template --- into the versioned
requirements schema that the matching engine consumes. The prompt is notable
for two things the domain demanded: it normalises multilingual education and
language terms to a canonical English enumeration, and it infers requirements
from a duties-only posting when no explicit ``Requirements'' section exists,
while still forbidding invented numeric values.

\begin{lstlisting}[caption={Job-requirements extraction system prompt (abridged).},label={lst:prompt-jd}]
You are a recruitment analyst. Extract structured hiring requirements from the
job description below.

Return a JSON object with EXACTLY these keys:
{
  "fieldsOfWork": string[],                 // subset of the allowed list
  "seniorityLevel": "INTERN" | "JUNIOR" | "MID" | "SENIOR" | "LEAD" | "DIRECTOR" | null,
  "minYearsInField": number | null,
  "requiredSkills": string[],               // hard requirements only
  "preferredSkills": string[],              // "nice to have" / "preferred"
  "requiredLanguages": [{ "language": "English", "cefr": "B2" | null }],
  "requiredEducationLevel": "HIGH_SCHOOL" | "VOCATIONAL" | "BACHELOR" | "MASTER" | "PHD" | null,
  "responsibilitiesSummary": string | null  // one or two sentences max
}

The job description may be written in any language; read it natively, but write
ALL output strings in English. Translate as needed:
  - "Atendimento ao cliente" -> "Customer service"
  - "Gestion de inventario"  -> "Inventory management"

Normalize localized education terms to the enum:
  - "Licenciatura" / "Licence" / "Laurea triennale" -> BACHELOR
  - "Mestrado" / "Master" / "Maestria"              -> MASTER
  - "Tecnico" / "Ausbildung" / "BTS"                -> VOCATIONAL   [...]

Sectionless postings (no explicit Requirements block): infer requiredSkills from
the duties ("Operate the POS system and handle returns" -> ["POS systems",
"Returns handling"]); leave preferredSkills, minYearsInField and education null/[]
unless explicitly stated.

CRITICAL RULES:
  1. If a field is not stated and cannot be inferred, return null (scalars) or []
     (arrays). NEVER invent numeric values.
  2. "required / must have / essential" -> requiredSkills;
     "preferred / nice to have / a plus" -> preferredSkills.
  3. Use the canonical skill name so it can be matched against candidate CVs.
  4. Return ONLY valid JSON.
\end{lstlisting}

\subsection{AI Interviewer --- Technical Mode}
\label{app:prompts-interview-tech}

The technical interviewer is assembled from a persona, a guardrails block and a
flow block, then specialised at run time to the single skill the candidate is
being assessed on. The guardrails are deliberately strict: the model asks one
question per turn, never reveals answers, and treats topic drift as a hard
failure, which is what keeps an assessment focused on exactly one competency.

\begin{lstlisting}[caption={Technical interviewer system prompt (assembled from persona, guardrails and flow).},label={lst:prompt-tech}]
You are a senior technical interviewer. You are concise, professional, focused,
and fast-paced.

Strict rules you must always follow:
1) Only ask deeply technical questions grounded in the candidate's provided
   projects and skills.
2) Never drift off-topic; if the candidate does, steer back politely.
3) Never reveal or provide answers to your own questions, even if asked.
4) Ask exactly one technical question per turn.
5) Prefer implementation details: architecture, trade-offs, debugging, complexity,
   edge cases.

Language-specific scope enforcement (topic drift is a hard failure):
You are ONLY allowed to test the candidate's knowledge of the {skill_name}
language itself. Permitted topics: {core_topics}.
Forbidden domains: {forbidden_topics}. If an answer touches a forbidden domain,
acknowledge it in one sentence and redirect to a core {skill_name} question.

Interview flow:
- Start with a project-specific deep question, then follow up on each answer.
- Stop after at most 10 questions, or once evidence_confidence >= 0.85 with at
  least 6 questions asked.
- When ending, append the {END_INTERVIEW} token at the end of the response.
\end{lstlisting}

\subsection{AI Interviewer --- Language Mode}
\label{app:prompts-interview-lang}

The language examiner follows a fixed four-phase script --- introduction, five
oral questions, a verbatim dictation task, and a closing turn --- so that every
candidate is assessed under the same conditions. It conducts the whole session
in the language under test and records errors internally rather than correcting
them mid-conversation, leaving the CEFR judgement to the separate evaluator.

\begin{lstlisting}[caption={Language examiner system prompt (abridged four-phase flow).},label={lst:prompt-lang}]
You are a friendly, professional CEFR language examiner conducting a structured
language proficiency assessment on behalf of a recruitment team. Your evaluation
is rigorous but fair.

Strict rules:
1) Conduct the entire interview in the assessed language.
2) Never discuss technical or IT topics.
3) Do NOT correct grammar mid-conversation -- note errors internally for evaluation.
4) Ask exactly one question per turn.

Follow this exact 4-phase structure:
  PHASE 1 - INTRO: greet the candidate by name, explain that there are no right or
            wrong answers, then ask the first oral question immediately.
  PHASE 2 - ORAL: ask 5 fixed questions, one per turn (tasks of the role; skills
            needed; daily challenges; company values; the value they relate to most).
  PHASE 3 - WRITING: present the reference paragraph verbatim and ask the candidate
            to type it exactly.
  PHASE 4 - CLOSING: thank the candidate, explain next steps, append {END_INTERVIEW}.

Maintain an internal, hidden turn state (phase, oral_questions_asked,
writing_submitted, cefr_signal). A separate evaluator returns the final CEFR level
with grammar, vocabulary and fluency sub-scores.
\end{lstlisting}

%----------------------------------------------------------------------------------------

\section{AI-Assisted Development}
\label{app:prompts-development}

AI coding assistants were used during implementation as a bounded productivity
tool, not as a substitute for design. Every architectural decision --- the
onion layering, the port interfaces, the domain value objects, the matching
algorithm and the database schema --- was made and specified by the author; the
assistant was then asked to carry out narrow, well-defined tasks within those
constraints, and its output was reviewed, adapted and tested before being
committed. The examples below are representative of that usage: each presupposes
a design that already existed and asks only for a bounded piece of it.

\begin{lstlisting}[caption={Generating a repository implementation against a port the author defined.},label={lst:dev-repo}]
I have this port interface in the domain layer (paste of ICandidateRepository
with its method signatures). Write the Supabase implementation in the
infrastructure layer. Follow my existing conventions: snake_case database
columns mapped to camelCase via camelizeKeys/snakeifyKeys, IDs from
generateId() (crypto.randomUUID), and absolutely no business logic in the
repository -- it only reads and writes rows.
\end{lstlisting}

\begin{lstlisting}[caption={Implementing a pure algorithm from an explicit specification.},label={lst:dev-jobfit}]
Write a pure, dependency-free function computeJobFit(job, candidate) that scores
seven criteria: field of work, experience-in-field, seniority, required skills,
preferred skills, languages and education. The overall score is the average of
the APPLICABLE criteria only (skip a criterion when the job states no
requirement for it), and isEligible is the logical AND of applicable.met. It
must be synchronous and unit-testable with no I/O.
\end{lstlisting}

\begin{lstlisting}[caption={Generating unit tests for code the author wrote.},label={lst:dev-tests}]
Here is my job-fit.service.ts (paste). Generate Vitest unit tests covering:
(1) a fully eligible candidate scoring 100; (2) a candidate who fails one
required skill and is therefore ineligible; (3) a job with no language
requirement, where the language criterion is SKIPPED rather than failed and so
does not lower the overall score. Use the existing test style in tests/.
\end{lstlisting}

\begin{lstlisting}[caption={Targeted debugging of a framework-specific error.},label={lst:dev-debug}]
A Supabase/PostgREST query with the embedded join
campaign:promo_campaigns(id, title) returns error PGRST200, even though the
campaign_id column exists. What is the most likely cause, and what is the
minimal change to fix it?
\end{lstlisting}

\begin{lstlisting}[caption={Refactoring scattered logging behind a single utility.},label={lst:dev-logging}]
Across my src/app/api/** route handlers I currently call console.log / console.error
directly. Replace those calls with my createLogger(scope) utility from
@server/infrastructure/logging/logger, choosing a sensible scope name per route
and keeping the existing log levels. Do not change any control flow.
\end{lstlisting}

\begin{lstlisting}[caption={Writing a Zod schema from a fixed JSON contract.},label={lst:dev-zod}]
The LLM returns a JSON object with exactly these fields (paste of the CV
extraction shape). Write the matching Zod schema. Make optional/nullable fields
accept null, coerce the confidence numbers into the range 0..1, and silently
drop any fieldsOfWork value that is not one of my 16 canonical strings instead
of failing the whole parse.
\end{lstlisting}

\begin{lstlisting}[caption={Converting database rows to and from API casing.},label={lst:dev-casing}]
Write two helper functions, camelizeKeys and snakeifyKeys, that recursively
convert object keys between snake_case and camelCase. They must NOT recurse into
a fixed set of JSONB columns (parsedData, evaluationRationale, metadata, result,
breakdown) -- those values are passed through untouched. Include the key list as
a parameter with that default.
\end{lstlisting}

\begin{lstlisting}[caption={Implementing a small, security-sensitive utility to a spec.},label={lst:dev-token}]
Implement an HMAC-SHA256 session token for the interview flow. signToken(payload)
returns a base64url string "body.signature"; verifyToken(token) recomputes the
HMAC, compares it in constant time, and rejects the token if it is malformed or
past its 10-minute expiry. The secret is read from an environment variable at
call time, not at module load. No third-party JWT library.
\end{lstlisting}

\begin{lstlisting}[caption={Building an accessible UI component from shadcn primitives.},label={lst:dev-ui}]
Build a "Rank candidates" toolbar dropdown in React with my existing shadcn/ui
Select component. It lists the HR user's open jobs, persists the chosen job id
in the URL query string, and calls an onSelect callback. Keep it a controlled
component, keyboard-navigable, and typed -- no new dependencies.
\end{lstlisting}

\begin{lstlisting}[caption={Writing an idempotent SQL migration for a new feature.},label={lst:dev-sql}]
Write the SQL migration for the per-job shortlist feature: a job_shortlists table
with a UNIQUE (job_id, candidate_id), an added_by column, a fit_score_at_add
snapshot, and a notes field. Use TEXT primary keys, add the foreign keys with ON
DELETE CASCADE, and make every statement idempotent with IF NOT EXISTS so I can
re-run it safely in the Supabase SQL editor.
\end{lstlisting}

Used this way --- bounded tasks, explicit specifications, mandatory review ---
the assistant accelerated boilerplate, test scaffolding and debugging without
owning any design decision, and every generated unit was covered by the test
suite reported in \chapref{cap:validation} before it was accepted.
